## chapter_01.tex

### 场景1：茶馆逼供

**场景编号**：1  
**地点**：鹅城茶馆，午后阳光透过破窗  
**人物**：六子、胡万、茶馆老板、围观群众  
**场景目的**：建立冲突，展示六子的困境  
**核心动作**：胡万把凉粉碗摔在地上，指着六子鼻子骂"吃了两碗只给一碗钱，穷鬼！"  
**冲突点**：六子掀桌子，碗碎声震耳欲聋，"老子没吃两碗！"  
**情感转变**：从屈辱愤怒 → 决绝爆发  
**场景结果**：胡万拔枪，六子被逼到墙角，生死一线

**正文**：

鹅城的午后像个蒸笼，茶馆的破窗户透进几缕昏黄的阳光，照在满是油污的八仙桌上。六子坐在角落里，面前摆着一碗已经凉透的凉粉，碗边还凝着一层白霜。

"吃啊，怎么不吃了？"胡万站在桌边，手里把玩着一把驳壳枪，枪套上的铜扣闪着冷光。他今天穿了件黑绸长衫，领口敞着，露出里面的白汗衫，一脸的横肉随着冷笑抖动。

六子的手在发抖，不是因为怕，是因为气。他明明只吃了一碗凉粉，给了一碗的钱，可胡万非说他吃了两碗。

"我就吃了一碗！"六子咬着牙，声音从牙缝里挤出来。

"一碗？"胡万突然暴喝，一脚踹翻了桌子。凉粉碗"啪"地一声摔在地上，碎成几瓣，凉粉溅了一地，混着尘土和碎瓷片，像摊烂泥。

茶馆里瞬间安静了，连苍蝇都不敢飞。老板缩在柜台后面，头都不敢抬。

"大家都看见了！"胡万指着地上的碎碗，又指着六子的鼻子，"这穷鬼吃了两碗凉粉，只给一碗钱！咱们鹅城的规矩，吃霸王餐要剁手！"

"我没吃两碗！"六子猛地站起来，椅子被他撞得向后倒去，发出刺耳的摩擦声。他的脸涨得通红，脖子上的青筋像蚯蚓一样暴起。

"没吃两碗？"胡万冷笑着，慢慢把手按在枪套上，"那这地上的碗是谁摔的？啊？"

六子的拳头攥得咯咯响，指甲掐进肉里。他看着胡万那张令人作呕的脸，看着周围人或是畏惧或是麻木的眼神，一股热血直冲脑门。

"老子没吃两碗！"六子吼出这句话的同时，双手猛地掀起桌子。实木的八仙桌在他手里像纸片一样飞起来，碗碟茶壶稀里哗啦砸了一地。

胡万没想到这小子敢动手，愣了一下，随即拔出枪，黑洞洞的枪口指着六子的眉心。

"找死！"胡万的手指已经搭在扳机上。

六子站在碎瓷片中，胸口剧烈起伏，眼神里没有恐惧，只有一种要把眼前这个人撕碎的狠劲。

茶馆里的空气凝固了，所有人都屏住呼吸，等着那声枪响。

### 场景2：绝地反杀

**场景编号**：2  
**地点**：茶馆同一位置，紧张对峙  
**人物**：六子、胡万、张麻子（暗处）  
**场景目的**：展示六子的反杀，引入张麻子视角  
**核心动作**：六子突然从怀里掏出左轮手枪，枪口直指胡万眉心  
**冲突点**：胡万冷笑"你敢开枪？"手指扣在扳机上  
**情感转变**：从绝望 → 绝地反击的快感  
**场景结果**：枪响，胡万倒地，全场死寂

**正文**：

时间仿佛凝固了。胡万的枪口离六子的眉心只有三寸，六子甚至能看见枪管里的来复线。

"跪下，磕头，叫三声爷爷，我就饶你一条狗命。"胡万的声音里带着猫戏老鼠的戏谑。

六子的腿在发抖，但他没有跪。他的手慢慢伸向怀里，动作很慢，像是在摸什么珍贵的东西。

胡万的眼睛眯了起来，手指在扳机上加了一分力。"别动！"

六子的手停在怀里，脸上突然露出一个诡异的笑。那笑容里没有恐惧，只有一种让人毛骨悚然的冷静。

"你知道西部牛仔是怎么决斗的吗？"六子轻声说。

胡万愣了一下："什么？"

就在这一瞬间，六子的手猛地抽出来——

"砰！"

一声巨响震得茶馆的窗户都在颤抖。

胡万的眉心多了一个血洞，他的表情还停留在错愕的那一刻，眼睛瞪得老大，似乎不敢相信发生了什么。然后，像截木头一样，他直挺挺地向后倒去，重重砸在地上，激起一片尘土。

枪是左轮手枪，枪管还冒着青烟。六子保持着射击的姿势，手腕稳如磐石，眼神冷得像冰。

茶馆里死一般的寂静。老板手里的算盘"啪"地掉在地上，围观的群众张大了嘴，连呼吸都忘了。

六子慢慢放下枪，枪口还指着胡万的尸体。他的手在微微颤抖，不是因为害怕，是因为后怕——如果慢了零点一秒，现在躺在地上的就是他。

"还有谁？"六子的声音不大，但在寂静的茶馆里像炸雷一样。

没有人说话。连苍蝇都不敢飞。

### 场景3：暗流涌动

**场景编号**：3  
**地点**：茶馆，枪响后的寂静  
**人物**：六子、黄四郎替身、张麻子、群众  
**场景目的**：引入黄四郎势力，张麻子认出西部枪法  
**核心动作**：六子用枪指着黄四郎替身："让真的黄四郎出来！"  
**冲突点**：替身冷笑"杀了胡万，你也活不过今晚"  
**情感转变**：从胜利的快感 → 意识到更大危机的恐惧  
**场景结果**：张麻子在暗处眯起眼睛，手摸向腰间的枪，认出那是牛仔的拔枪术

**正文**：

"好枪法。"

一个阴冷的声音从茶馆二楼传来。

六子猛地抬头，看见一个穿着和胡万一模一样黑绸长衫的人站在楼梯上，手里也把玩着一把驳壳枪。那人的脸和胡万有七分像，但眼神更阴鸷，嘴角挂着一丝让人不舒服的笑。

这是黄四郎的替身。鹅城人都知道，黄四郎从不以真面目示人，他有无数个替身，没人知道哪个是真的。

"你是谁？"六子的枪口转向楼梯上的人。

"我是谁不重要。"替身慢慢走下楼梯，皮鞋踩在木地板上，发出有节奏的"咚咚"声，"重要的是，你杀了胡万，就是打了黄老爷的脸。"

茶馆里的群众开始往后退，有人已经悄悄溜出门去。他们知道，接下来要发生的事，不是他们能看的。

"让真的黄四郎出来！"六子吼道，枪口稳稳指着替身的眉心。

替身停在离六子五步远的地方，冷笑一声："就凭你？一个吃凉粉都给不起钱的穷鬼？"

"我不是穷鬼！"六子的手指扣在扳机上，"我是张麻子的儿子！"

这句话一出，茶馆里响起一片倒吸冷气的声音。张麻子，鹅城最近出现的神秘人物，传说中的土匪，没人见过他的真面目，但人人都怕他的名字。

替身的眼神变了变，但很快恢复了阴冷："张麻子？哼，他来了也得给黄老爷跪下。杀了胡万，你活不过今晚。"

"那就试试！"六子向前迈了一步，枪口几乎顶到替身的鼻子上。

替身的手慢慢按在枪套上，两人的距离只有三步，这是生死的距离。

就在这时，茶馆角落的阴影里，一双眼睛正冷冷地注视着这一切。

张麻子坐在最暗的角落，宽檐帽压得很低，只露出下半张脸。他的手指轻轻摩挲着腰间的 Colt Peacemaker 左轮手枪，那是他从西部带来的老伙计。

刚才六子拔枪的那一瞬间，张麻子的眼睛亮了。

那是标准的西部拔枪术——手腕不动，靠小臂的爆发力瞬间出枪，枪口直指眉心，没有一丝多余的动作。这种枪法，在民国的鹅城，除了他，不应该有第二个人会。

"这小子..."张麻子低声自语，嘴角勾起一丝不易察觉的笑，"有点意思。"

他慢慢站起来，阴影吞没了他的身影，只留下腰间的枪套在阳光下闪着冷光。

茶馆里的对峙还在继续，但张麻子知道，这只是开始。

真正的戏，才刚刚开场。